\documentclass{jsarticle}
\usepackage[dvipdfmx]{graphicx}
\usepackage{amsmath}
\usepackage{amssymb}
\begin{document}
\title{Deprivate Equation, Global assemble equation}
\author{Masaaki Yamaguchi}
\date{}
\maketitle


   大域的多重積分と大域的無限微分多様体を定義すると、

   $$ \int f(x)dx = \int \Gamma(\gamma)^{'}dx_m $$
   $$ = 2 (\cos (i x \log x) - i \sin(i x \log x)) $$
   $$ \left({{\int f(x) dx} \over \log x}\right) = \lim_{\theta \to \infty}
    \left({{\int f(x) dx} \over \theta}\right) = 0,1 $$
    $$ e^{i \theta} = \cos \theta + i \sin \theta $$
    $$ \left({\int f(x) dx}\right)^{'} = 2 (i \sin(i x \log x) - \cos (i x \log
    x)) $$
    $$ = 2(- \cos(i x \log x) + i \sin(i x \log x)) $$
    $$ (\cos (i x \log x) - i \sin(i x \log x))^{'} $$
    $$ = {d \over {d {e^{i\theta}}}}\left((\cos, - \sin )\cdot
    (\sin, \cos)\right)  $$
    $$ \int \Gamma(\gamma)^{''}dx_m = \left({\int \Gamma(\gamma)^{'}}dx_m
    \right)^{\nabla L} = \left({{\int \Gamma dx_m} \cdot {{d \over d\gamma}
    \Gamma}}\right)^{\nabla L} \le 
    \left({\int \Gamma dx_m + {d \over d\gamma}\Gamma}\right)^{\nabla L} \le 
    \left({e^{f} - e^{-f}} \le {e^{-f} + e^{f}}\right)^{'}  $$
    $$ = 0,1 $$

    以上より、大域的微分多様体を大域的2重微分多様体として、
    処理すると、ホモロジー多様体では、種数が１であり、特異点では、種数が０
    と計算されることになる。ガンマ関数の大域的微分多様体では、
    シュバルツシルト半径として計算されうるが、これを大域的2重微分で
    処理すると、ブラックホールの特異点としての解が無になる。

    川島隆太先生が、私みたいな人が、頭の中で偏微分や多重積分、部分積分、
    共変微分や、置換積分や、複素微分や複素積分を考えていると思っていた
    と言っていたが、それ以上の大域的微分や大域的2重微分や、
    大域的多重積分や大域的偏微分や大域的部分積分、微分幾何の量子化、
    単体積分、単体微分などの、大学院生でも考えない、思いもしない、
    実際の式を書くと誰でも、その関数と多様体が存在するなと思うしか
    できない、私の記述式を見ると納得する理論式でもあると、
    私も自慢できる理論であると自負できる。

  $$ \log x|_{g_{ij}}^{\nabla L} = f^{f^{'}}, F^f|_{g_{ij}}^{\nabla L}
  |: x \to y, x^p \to y $$ 
  $$  f(x) = \log x = p \log x, f(y) = p \log x  $$
  大域的偏微分方程式と大域的多重積分は、それぞれ次のように成り立っている。
  $$ {d^2 \over df^2}F = F^{f^{'}}\cdot f^{f^{''}},  $$
  $$ \int \int F dx_m = F^f \cdot F^{(f)^{'}} $$
  $$ {d \over df dg}(f,g) = (f \cdot g)^{f^{'} + g^{'}} $$
  大域的部分積分も、次のように成り立っている。
  $$ (F^f \cdot G^g) = \int (f \cdot g)^{f^{'} + g^{'}} $$
  $$ \int \int F\cdot G dx_m = [F^f \cdot G^g] - 
  \int (f \cdot g)^{f^{'} + g^{'}} $$
  大域的部分積分の計算は、
  $$ \int {d \over {df dg}}FG = \int F^{f^{'}}G + \int F G^{g^{'}} $$
  $$ \int F^{f^{'}}Gdx_m = [F^f G^g] - \int F G^{g^{'}}dx_m$$
  大域的商代数の計算は、
  $$ \left({F(x) \over G(x)}\right)^{(fg)^{'}} = {F^{f^{'}}G - FG^{g^{'}}
  \over G^{g}} $$
  大域的偏微分方程式は、縮約記号を使うと、
  $$ {d \over df dg}FG = {d \over df_m}FG $$
  $$ {\partial \over \partial f_m}FG = F^{f^{\mu\nu}}\cdot G
  + F \cdot G^{g^{\mu\nu}}, \int F dx_m = F^f $$
  多様体による大域的微分と大域的積分が、エントロピー式で統一的に表せられる。

   ベータ関数をガンマ関数へと渡すと、
   $$ \beta(p,q) = {{\Gamma(p)\Gamma(q)} \over {\Gamma(p + q)}} $$
   ここで、単体的微分と単体的積分を定義すると
   $$ t = \Gamma(x) $$
   $$ T = \int \Gamma(x) dx $$
   これを使うのに、ベータ関数を単体的微分へと渡すために、
   $$ \beta(p,q) = - \int {1 \over {t^2}}dt $$
   大域的微分変数へとするが、
   $$ T_m = \int \Gamma(x) dx_m $$
   この単体的積分を常微分へとやり直すと、
   $$ T^{'} = {t^{'} \over {{\log t}}}dt + C(\text{Cは積分定数}) $$
   これが、大域的微分多様体の微分変数と証明するためには、
   $$ dx_m = {1 \over {\log x}}dx, dx_m = {(\log x^{-1})}^{'}  $$
   これより、単体的微分が大域的部分積分へと置換できて、
   $$ T^{'} = \int \Gamma(\gamma)^{'} dx_m $$
   微分幾何へと大域的積分と大域的微分多様体が、加群分解の共変微分へと
   書き換えられる。
   $$ \bigoplus T^{\nabla}  = \int T dx_m, \delta(t) = t dx_m  $$


   これらをまとめると、ベータ関数を単体的関数として、
   これを常微分へと解を導くと、単体的微分と単体的積分
   へと定義できて、これより、大域的微分多様体と大域的積分多様体
   へとこの関数が存在していることを証明できるのを上の式たちから
   わかる。この単体的微分と単体的積分から微分幾何が書けることがわかる。



   虚数の虚数倍した値が超越数のx倍と同じとすると、
   超越数の$2 \pi m$倍が$n=1$でiとなると定義すると、
   次の式たちが導かれる。　


   カルタンを超えているアイデアと数式たちでもある。

   $$ i^i = (\sqrt{i})^{\sqrt{i}} = e^{x \log x} $$
   $$ e^x = i, e^{2\pi m} = i^n, {d \over d x}e^{2\pi m} = i^n $$
   $$ e^{i \theta} = \cos \theta + i \sin \theta, e^{2 \pi m} = i^n $$
   $$ 2 \pi m = n, e = i \text{($e=i$となる$e^{i \theta} = \cos \theta 
   + i \sin \theta $
   の範囲で$2\pi m$ がある。)}$$
   $$ {d \over d i}i = i^i = e^{x \log x}, m={1 \over 2 \pi},l=2\pi r, 
   \pi = {l \over 2r} $$
   $$ ||ds^2|| = e^{-2 \pi T||\psi||}[\eta_{\mu\nu} + \bar{h}(x)]
   dx^{\mu}dx^{\nu} + T^2 d^2 \psi $$
   $$ e^x = i, e^{2\pi m} = i^n, e = i $$
   $$ {{[\eta_{\mu\nu} + \bar{h}(x)]dx^{\mu}dx^{\nu}}/i} = H \Psi
   = i \hbar \psi, H \Psi = {1 \over i}[H, \Psi] $$
   ハイゼンベルク方程式が$ AdS_5 $多様体の原子レベルの方程式も表せられる。 
   微分幾何の量子化の式は、
   $$ \bigoplus {({i \hbar} ^{\nabla})}^{\oplus L} = e^{{i x \log x}^{
	   {e^{i x \log x}}^{'}}} $$
   $$ \Gamma = \int e^{-x}x^{1 - t}dx $$
   $$ \gamma = \int e^{-x}x^{1 - t}\log x dx $$
   $$ = \int \Gamma(\gamma)^{'} dx_m$$
   $$ = {d \over d \gamma}\Gamma $$

   ８種類の微分幾何では、それぞれの時間の固有値が違う。
   閉3次元多様体上で微分幾何の切除、分解によって時間の性質が決まる。
   閉３次元多様体に統合されると、ゼータ関数になり、各幾何に分解された場合
   に、非分岐から、この上、sergeryの結果が決まる。

   各微分構造においての時間発展での熱エネルギーの変化は$E=mc^2, E=m^2 c^2 $
   のthree manifoldが微分幾何構造の変化にそれぞれ対応する。

   $$ \nabla ({i \hbar^{\nabla}})^{\oplus L}, \bigoplus ({i \hbar^{\nabla}})^{
	   \oplus L}, \square ({i \hbar^{\nabla}})^{\oplus L} $$
   これらの作用素が加減乗除の生成式の元、Euler Productの結論による
   作用素生成の論理素子でもある。
   その上に、8種類の微分幾何は、Jones多項式の4パターンでの構成される
   差分と加群方程式から、特殊相対性理論をも思わせる、この差分方程式
   $$ {d \over d \gamma}\Gamma = e^{f} + e^{-f} \ge e^{f} - e^{f} $$
   から、8種類の代数幾何が、ストークスの定理と同じく、このJones多項式
   の固有周期が、すべての時間の流れの基軸をも内包してもいることが、
   わかる。このJones多項式は、結び目の理論をも、この周期が言い表してもいる。
   $$ {d \over d \gamma}\Gamma + \int C dx_m = 2(\cos(i x \log x) - i \sin
   (i x \log x)) $$
    $$ e^{- \theta} = \cos(i x \log x) - i \sin(i x \log x)$$
    $$ \int \Gamma(\gamma)^{'}dx_m = 2(\cos (i x \log x) 
    - i \sin(i x \log x)) $$

  ８種類の微分幾何での、それぞれの時間の固有値の動静は、
  ストークスの流体力学での、熱エネルギーの各流れの仕方で表されている。
  種数１の2種が種数０の球体に両辺で交わっていると、
一般相対性理論における重力理論は大域的微分多様体と積分多様体についての 単体量を求めるための空間の対数関数における不変性を記述するプランクスケール と異次元の宇宙におけるスケール、ウィークスケールと言われているanti-D-brane として、この２種類の計量をガンマ関数とベータ関数としてオイラーの定数 とそれによる連分数が微分幾何の量子化と数式の値を求めると同じという 予知と推測値から、広中平祐定理の４重帰納法のオイラーの公式による 多様体積分と同類値として、ペレルマン多様体がサーストン空間に成り立つ と同じく、広中平祐定理がヒルベルト空間で成り立つとしての、 これら２種の多様体の場の理論が、ゼータ関数がAdS5多様体で成り立つことと、 不変量として、ゼータ関数がこの3種の多様体の場の理論のバランスをとる
  理論として言えることである。
$$ ||ds^2|| = \lim_{x \to \infty} [\delta(x) \int \int \int \pi \left( \sum_{k=0}^{\infty}
		{{}^n\sqrt{p},x \over n} \right)
^{1 \over 2}d \tau]^{\mu\nu} $$
  This norm is component of fields in Hilbert manifold of space theorem rebuilt with
Mebius space in Gamma Function on Beta function fill into power of rout fundamental
group.

  　 Gravity of general relativity theory describe with Cutting of space
in being discatastrophed from Global differential and integral manifold
of scaled levetivity in plank scal and the other vector of universe
scale, these two scale inspectivity of Gamma and Beta function escourted
with these manifold experted in result of Differentail geometry of
quantum level manifold equal with Euler product and continue parameter,
moreover this product is relativity of Hironaka theorem in Four assembled
integral of Euler equation.

   \newcommand{\variint}{%
	\begin{picture}(15,30)
		\put(0,0){
			$ \iint $}
		\put(0,5){\line(15,0){15}}
	\end{picture} %
}

\newcommand{\variiint}{%
	\begin{picture}(15,15)
		\put(0,0){
			$ \iiint $}
		\put(0,5){\line(15,0){20}}
	\end{picture} %
}



 \newcommand{\alearletter}{%
	 \begin{picture}(10,20)
		 \put(0,0){
			 $ \square $}
		 \put(0,0){\line(1,1){10}}
	 \end{picture} %
	 }

 重力場理論の式は、
  Gravity equation is
  $$ \square = - {16 \pi G \over c^4}T^{\mu\nu} $$
  This equation quated with being logment of formula, and
  this formula divided with universe of number in prime zone,
  therefore, this dicided with varint equation is monotonicity of
  being composited with Weil's theorem united for Gamma function.
  この式は、対数を宇宙における数により求める素数分布論として、この
  大域的積分分断多様体がガンマ関数をヴェイユ予想を根幹とする単体量
  として決まることに起因する商代数として導かれる。
  $$ \scalebox{1}[2]{\alearletter} = 
  {{{{8 \pi G} \over {c^4}}T^{\mu\nu}}/ \log x} $$
  $$ {}^{t}\scalebox{1}[2]{\variiint}\mathrm{cohom D_{\chi}[I_m]}$$
  $$ = \oint {(px^n + qx + r)^{\nabla l}} $$
  $$ {d \over d l}L(x,y) = 2 \int ||\sin 2x||^{2}d \tau $$
  $$ {d \over d \gamma}\Gamma $$
  この関数は大域的微分多様体としてのアカシックレコードの合流地点として、
  タプルスペースを形成している。
  This function esterminate with acasic record of global differential
  manifolds.
  $$ = [i \pi (\chi, x), f(x)] $$
  それにより、この多様体は基本群をアカシックレコードの相対性としての
  存在論の実存主義から統合される多様体自身としてのタプルスペースの
  池になっている。
  And this manifold from fundemental group esterminate with
  also this manifold estimate relativity of acasic record.
  $$ {d \over d \gamma}\Gamma = \int C dx_m = \int(\int {1 \over x^s}dx
  - \log x)\mathrm{dvol} $$
  また、このアカシックレコードはオイラーの定数のラムダドライバー
  にもなっている。
  More also this record tupled with lake of Euler product.
  $$ \bigoplus({i \hbar^{\nabla}})^{\oplus L} $$
  $$ x^{{1 \over 2} + iy} = e^{x \log x} $$
  それゆえに、この定数はゼータ関数と微分幾何の量子化を因数にもつ
  素因子分解の式にもなっている。
  Therefore, this product also constructed with differential geometry
  of quantum level equation and zeta function.
  $$ = C $$ 
  そして、この関数はオイラーの定数から広中平祐定理による4重帰納法の
  オイラーの公式からの多様体積分へとのサーストン空間のスペクトラム関数
  ともなっている。
  And this function of Euler product respectrum of focus with
  Heisuke Hironaka manifold in four assembled of integral Euler equation.. 

  $$ C =  b_0 + \cfrac{c_1}{b_1 +
	  \cfrac{c_2}{b_2 +
	  \cfrac{c_3}{b_3 +
	  \cfrac{c_4}{b_4 + \cdots}}}} $$
  この方程式は指数による連分数としての役割も担っている。  
  This equation demanded with continued number of step function.
  $$ x^{{1 \over 2} + iy} = e^{x \log x} $$
  $$ (2.71828)^{2.828} = a^{a + b^{b + c^{c + d^{d \cdots }}}} $$
  $$ = \int e^f \cdot x^{1 - t}dx $$
  This represent is Gamma function in Euler product.
  Therefore this product is zeta function of global differential equation.

  $$ {d \over d f}F(v_{ij},h) = \int {e^{-f}\mathrm{dV}}
  [- \Delta v + \nabla_{i}\nabla_{j}
  v_{ij} - R_{ij}v_{ij} - v_{ij}\nabla_{i}\nabla_{j} + 2 <\nabla f,\nabla h>
  + (R + \nabla f^2)({v \over 2} - h)]$$
  $$ = {d \over d f}F = {2 {\int (R + \nabla_{i}\nabla_{j} f)^2} \over 
  { - (R + \Delta f)}}\mathrm{dm} $$
  これらの方程式は8種類の微分幾何の次元多様体として、そして、これらの
  多様体は曲平面による双対性をも生成している。 
  そして、このガウスの曲平面は、大域的微分多様体と微分幾何の量子化から
  素因数を形成してもいる。
  These equation are eight differential geometry of dimension calyement.
  And these calyement equation excluded into pair of dimension surface.
  This surface of Gauss function are global differential manifold,
  and differential geometry of quantum level.
  $$ F \ge {d \over d f}\int \int{1 \over (x \log x)^2}dx_m 
  + {d \over d f}\int \int{1 \over (y \log y)^{1 \over 2}}dy_m $$

  微分幾何の量子化はオイラーの定数とガンマ関数が指数による連分数
  としての不変性として素因数を形成していて、このガウスの曲平面
  による量子力学における重力場理論は、ダランベルシアンの切断多様体
  がこの大域的切断多様体を付加してもいる。
  Differential geometry of quantum level constructed with Euler product
  and Gamma function being discatastrophed from continued fraction style.
  Gravity equation lend with varint of monotonicity of level expresented
  from gravity of letter varient formula.
  これらの方程式は基本群と大域的微分多様体をエスコートしていもいて、
  ヴェイユ予想がこのダランベルシアンの切断方程式たちから
  輸送のポートにもなっている。ベータ関数とガンマ関数がこれらのフォームラ
  の方程式を放出してもいて、結果、これらの方程式は広中平祐定理の複素
  多様体とグリーシャ教授によるペレルマン多様体からサポートされてもいる。
  この2名の教授は、一つは抽象理論をもう一方は具象理論を説明としている。
  These equation escourted into Global differential manifold and
  fundemental equation.
  Weil's theorem is imported from this equation in gravity of letters.
  Beta function and Gamma function are excluded with these formulas.
  These equation comontius from Heisuke Hironaga of complex manifold
  and Gresha professor of Perelman manifold.
  These two professos are one of abstract theorem and the other of
  visual manifold theorem.

\section{Hilbert manifold in Mebius space \\
        this element of Zeta function on integrate of fields}

 Hilbert manifold equals with Von Neumann manifold, and this fields is concluded with Glassman manifold,
the manifold is dualty of twister into created of surface built.
These fields is used for relativity theorem and quantum physics.
$$ ||ds^2|| = \lim_{x \to \infty} [\delta(x) \int \int \int \pi \left( \sum_{k=0}^{\infty}
		{{}^n\sqrt{p},x \over n} \right)
^{1 \over 2}d \tau]^{\mu\nu} $$
  This norm is component of fields in Hilbert manifold of space theorem rebuilt with
Mebius space in Gamma Function on Beta function fill into power of rout fundamental
group.
And this result with $AdS_5$ space time in Quantum caos of Minkowsky of manifold in
abel manifold.
Gravity of metric in non-commutative equation is Global differential equation
conbuilt with Kaluza-Klein space. Therefore this mechanism is $T^{\mu\nu}$ tensor
is equal with $R^{\mu\nu}$ tensor. And This moreover inspect with laplace operator
in stimulate with sign of differential operator.
Minus of zone in Add position of manifold is Volume of laplace equation rebuilt with
Gamma function equal with summative of manifold in Global differential equation,
this result with setminus of zone of add summative of manifold, and construct with
locality theorem straight with fundamental group in world line of surface, this
power is boson and fermison of cone in hyper function.

$$ V(\tau) = [f(x),g(x)] \times [f^{-1}(x), h(x)] $$
$$ \Gamma(p,q) = \int e^{-x} x^{1 - t}dx $$
$$ = \beta(p,q) $$
$$ = \pi (f(\chi,x),x) $$
$$ ||ds^2|| = \mathcal{O}(x)[(f(x) \circ g(x))^{\mu\nu}]dx^{\mu}dx^{\nu} $$
$$ = \lim_{x \to \infty} \sum_{k=0}^{\infty}a_k f^k $$
$$ G^{\mu\nu} = {\partial \over \partial f}\int [f(x)^{\mu\nu} \circ G(x)^{\mu\nu}dx^{\mu}dx^{\nu}]^{\mu\nu}dm $$
$$ = g_{\mu\nu}(x)dx^{\mu}dx^{\nu} - f(x)^{\mu\nu}dx^{\mu}dx^{\nu} $$
$$ [i \pi(\chi,x),f(x)] = i \pi f(x) - f(x) \pi(\chi,x) $$
$$ T^{\mu\nu} = (\lim_{x \to \infty} \sum_{k=0}^{\infty}\int \int [V(\tau) \circ S^{\mu\nu}(\chi,x)]dm)
^{\mu\nu}dx^{\mu}dx^{\nu} $$
$$ G^{\mu\nu} = R^{\mu\nu}T^{\mu\nu} $$
$$ \begin{vmatrix} D^m & dx \\
dx & \sigma^m \end{vmatrix} \begin{vmatrix} \cos \theta & - \sin \theta \\
\sin \theta & \cos \theta \end{vmatrix} \begin{vmatrix} x \\
y \end{vmatrix} = \begin{pmatrix} 1 & 0 \\ 
0 & - 1 \end{pmatrix}^{1 \over 2} $$
$$ \sigma^m \begin{bmatrix} \delta(x) & - 1 \\
1 & \epsilon(x) \end{bmatrix}^{1 \over 2} = \begin{pmatrix} i & 0 \\
0 & - i \end{pmatrix} $$
$$ V(M) = {\partial \over \partial f} ({}^N \int [f \smallsetminus M]^{\oplus N})^{\mu\nu} dx^{\mu}dx^{\nu} $$
$$ V(M) = \pi(2 \int \sin^2 dx)\oplus {d \over df}F^M dx_m $$
$$ \lim_{x \to \infty}\sum_{k=0}^{\infty} a_k f^k = \int (F(V)dx_m)^{\mu\nu}dx^{\mu}dx^{\nu} $$
$$ \bigoplus_{k=0}^{\infty}[ f \smallsetminus g] = \vee(M \wedge N) $$
$$ \pi_1 (M) = e^{-f 2 \int \sin^2 x dm} + O(N^{-1}) $$
$$ = [i \pi(\chi,x),f(x)] $$
$$ M \circ f(x) = e^{-f \int \sin x \cos x dx_m} + \log (O(N^{-1})) $$

 Non-Symmetry space time.
 $$ {d \over dL}V(\tau) = {d \over df}\int \int_M {1 \over (x \log x)^2}dx_m + {d \over df}\int \int_M
 {1 \over (y \log y)^{1 \over 2}}dy_m $$
 $$ \epsilon S(\nu) = \square_v \cdot {\partial \over \partial \chi}({}^5 \sqrt{\wedge g^2})d\chi $$
 Differential Volume in $AdS_5$ graviton of fundamental rout of group.
 $$ \wedge (F^m_t)^{''} = {1 \over 12}g_{ij}^2$$
 Quarks of other dimension.
 $$ \pi(V_{\tau}) = {e^{-(\sqrt{\pi \over 16}\log x)}}^{\delta} \times {1 \over (x \log x)} $$
 Universe of rout, Volume in expanding space time.
 $$ {d \over dt}(g_{ij})^2 = {1 \over 24}(F^m_t)^2 $$
 $$ m^2 = 2 \pi T \left({26 - D_n \over 24}\right) $$
 This quarks of mass in relativity theorem, and fourth of universe in three manifold of one dimension surface,
 and also this integrate of dimension in conbuilt of quarks.
 $$ g_{ij} \wedge \pi(\nu_{\tau}) = e^{-2 \pi T |\psi|}[\eta_{\mu\nu} + \bar{h}_{\mu\nu}(x)]dx^{\mu}dx^{\nu} 
 + T^2 d \psi^2 $$
 Out of rout in $AdS_5$ space time.
 $$ ||ds^2|| = g_{ij}\wedge \pi(\nu_v) $$
 $AdS_5$ norm is fourth of universe of power in three manifold out of rout.


  Sphere orbital cube is on the right of hartshorne conjecture by the right equation,
this integrate with fourth of piece on universe series, and this estimate with one of three manifold.
Also this is blackhole on category of symmetry of particles.
These equation conbuilt with planck volume and out of universe is phologram field, this field construct with
electric field and magnitic field stand with weak electric theorem.
This theorem is equals with abel manifold and $AdS_5$ space time.
Moreover this field is anti-brane and brane emerge with gravity and antigravity equation restructed with.
Zeta function is existed with Re pole of ${1 \over 2}$ constance reluctances.
This pole of constance is also existed with singularity of complex fields.
Von Neumann manifold of field is also this means.
$$ ||ds^2|| = \lim_{x \to \infty} [\delta(x)\int \int \int \pi \left( \sum_{k=0}^{\infty}{{}^n\sqrt{p},x \over n}\right)
^{1 \over 2}d \tau]^{\mu\nu} $$
$$ e^{i \theta} = \cos \theta + i \sin \theta $$
$$ e^{x \log x} = x^{{1 \over 2} + iy}, x \log x = \log (\cos \theta + i \sin \theta) $$
$$ = \log \cos \theta + i \log \sin \theta $$
This equation is developed with frobenius theorem activate with logment equation.
This theorem used to
$$ \log(\sin \theta + i \cos \theta)= \log(\sin \theta - i \cos \theta) $$
$$ \log \left({\sin \theta \over i \cos \theta}\right) = -2 R_{ij}, {d \over dt}g_{ij}(t) =  - 2 R_{ij} $$
$$ \mathcal{O}(x) = {\zeta(s) \over \sum\limits_{k=0}^{\infty}a_k f^k} $$
$$ \mathrm{Im} f = \mathrm{ker} f, \chi(x) = {\mathrm{ker}f \over \mathrm{Im}f} $$ $$ H(3) = 2, \nabla H(x) = 2, 
 \pi(x) = 0 $$
 This equation is world line, and this equation is non-integrate with relativity theorem of rout constance.

$$ [f(x)] = \infty, ||ds^2|| = \mathcal{O}(x)[\eta_{\mu\nu} + \bar{h}_{\mu\nu}(x)]dx^{\mu}dx^{\nu} + T^2 d^2 \psi $$
$$ T^2 d^2 \psi = [f(x)] , T^2 d^2 \psi = \lim_{x \to 1}\sum_{k=0}^{\infty}a_k f^k $$

 These equation is equals with M theorem. For this formula with zeta function into one of universe in four dimensions.
Sphere orbital cube is on the surface into $AdS_5$ space time, and this space time is Re pole and Im pole of 
constance reluctances.


  Gauss function is equals with Abel manifold and Seifert manifold. Moreover this function is infinite time,
and this energy is cover with finite of Abel manifold and Other dimension of seifert manifold on the surface.


  Dilaton and fifth dimension of seifert manifold emerge with quarks and Maxwell equation, this power
is from dimension to flow into energy of boson. 
This boson equals with Abel manifold and $AdS_5$ space time, and this space is created with Gauss function.


  Relativity theorem is composited with infinite on D-brane and finite on anti-D-brane, This space time
restructed with element of zeta function.
Between finite and infinite of dimension belong to space time system, estrald of space element have with 
infinite oneselves. Anti-D-brane have infinite themselves, and this comontend with fifth dimension of AdS5 
have for seifert manifold, this asperal manifold is non-move element of anti-D-brane and move element 
of D-brane satisfid with fifth dimension of seifert algebra liner.
Gauss function is own of this element in infinite element.
And also this function is abel manifold.
Infinite is coverd with finite dimension in fifth dimension of AdS5 space time.
Relativity theorem is this system of circustance nature equation.
AdS5 space time is out of time system and this system is belong to infinite mercy.
This hiercyent is endrol of memolite with geniue of element.
Genie have live of telomea endore in gravity accesorlity result.
AdS5 space time out of over this element begin with infinite assentance.
Every element acknowlege is imaginary equation before mass and spiritual envy.

  $$ T^2 d^2 \psi = \lim_{x \to 1}\sum_{k=0}^{\infty}a_k f^k $$
  $$ \lim_{x \to 1}\sum_{k=0}^{\infty}a_k f^k = [T^2 d^2 \psi] $$

  $$ {d \over dL}V(\tau)= {d \over df}\int \int_M({}^5\sqrt{x^2})d\Lambda 
  + {d \over df}\int \int_M {}^N({}^3\sqrt{x})^{\oplus N}d\Lambda $$
  $$ {}^M(\vee(\wedge f \circ g)^N)^{1 \over 2} = {d \over df}\int \int_M{1 \over (x \log x)^2}dx_m
  + {d \over df}\int \int_M{1 \over (y \log y)^{1 \over 2}}dy_m $$
  $$ ||ds^2|| = \mathcal{O}(x)[\eta_{\mu\nu} + \bar{h}_{\mu\nu}(x)]dx^{\mu}dx^{\nu}
  + T^2 d^2 \psi $$
  $$ \mathcal{O}(x) = e^{-2\pi T|\psi|} $$
  $$ G^{\mu\nu} = R_{\mu\nu} T^{\mu\nu} $$
  $$ = - {1 \over 2}\Lambda g_{ij}(x) + T^{\mu\nu} $$

  Fifth dimension of eight differential structure is integrate with one geometry element,
Four of universe also integrate to one universe, and this universe represented with symmetry formula.
Fifth universe also is represented with seifert manifold peices.
All after universe is constructed with six element of circuatance.
This aquire maniculate with quarks of being esperaled belong to.

\section{Imaginary equation in AdS5 space time create with \\
       dimension of symmetry}

  D-brane and anti-D-brane is composited with all of series universe emerged for one geometry of 
dimension, this gravity of power from D-brane and anti-D-brane emelite with ancestor.
Seifert manifold is on the ground of blackhole in whitehole of power pond of senseivility.
Six of element of quarks and universe of pieces is supersymmetry of mechanism resolved with

hyper symmetry of quarks constructed to emerge with darkmatter.
This darkmatter emerged with big-ban of heircyent in circumstance of phenomena.

  D-brane and anti-D-brane equations is
  $$ {d \over dL}V(\tau) = {d \over df}\int \int_M {1 \over (x \log x)^2}dx_m
  + {d \over df}\int \int_M {1 \over (y \log y)^{1 \over 2}}dy_m $$
  $$ C = \int \int {1 \over (x \log x)^2}dx_m $$

  Euler constance is quantum group theorem rebuild with projective space involved with.


   虚数方程式は、反重力に起因するフーリエ級数の励起を生成する。それは、人工知能を生み出す、
   ５次元時空にも、この虚数方程式は使われる。AdS5の次元空間は、反ド・シッター時空の
   D-braneとanti-D-braneのcomformal場を生み出す。ホログラフィー時空は、この量子起因によるものである。
   2次元曲面によるブラックホールは、ガンマ線バーストによる５次元時空の構造から観測される。
   空間の最小単位によるプランク定数は、宇宙の大域的微分多様体から導き出される、AdS5の次元空間の
   準同型写像を形成している。これは、最小単位から宇宙の大きさを導いている。最大最小の方程式は、
   相加相乗平均を形成している。時間と空間は、宇宙が生成したときから、宇宙の始まりと終わりを既に生み出している。
   宇宙と異次元から、ブラックホールとホワイトホールの力がわかり、反重力を見つけられる。オイラーの定数は、
   この量子定数からわかる。虚数の仕組みはこの量子スピンの産物である。
   オイラーの定数は、 この虚時空の斥力の現存である。それは、非対称性理論から導かれる。
   不確定性原理は、AdS5のブラックホールとホワイトホールを閉３次元多様体に統合する５次元時空から求められる。
   位置と運動エネルギーが、空間の最小単位であるプランク定数を宇宙全体にする微細構造定数からわかり、
   面積確定から、アーベル多様体を母関数に極限値として、ゼータ関数をこの母関数に不変式として、
   ２種類ずつにまとめる４種類の宇宙を形成する８種類のサーストンの幾何化予想から導き出される。
   この閉３次元多様体は、ミラー対称性を軸として、６種類の次元空間を一種類の異次元宇宙と同質ともしている。
   複素多様体による特異点解消理論は、この原理から求められる。この特異点解消理論は、２次元曲面を３次元多様体
   に展開していく、時空から生成される重力の密度を反重力と等しくしていく時間空間の４次元多様体と虚時空
   から求められる。
   ヒルベルト空間は、フォン・ノイマン多様体とグラスマン多様体をこのサーストンの幾何化予想を場の理論
   既定値として形成される。
   この空間は、ミンコフスキー時空とアーベル多様体全体を表している。そして、この空間は、球対称性を
   複素多様体を起点として、大域的トポロジーから、偏微分を作用素微分として時空間をカオスから
   ずらすと５次元多様体として成立している。
   これらより、３次元多様体に２次元射影空間が異次元空間として、AdS5空間を形成される。
   偏微分、全微分、線形微分、常微分、多重微分、部分積分、置換微分、大域的微分、単体分割、
   双対分割、同調、ホモロジー単体、コホモロジー単体、群論、基本群、複体、マイヤー・ビーートリス完全系列、
   ファン・カンペンの定理、層の理論、コホモルティズム単体、CW複体、ハウスドロフ空間、線形空間、位相空間、
   微分幾何構造、モースの定理、カタストロフの空間、ゼータ関数系列、球対称性理論、スピン幾何、
   ツイスター理論、双対被覆、多重連結空間、プランク定数、フォン・ノイマン多様体、グラスマン多様体、
   ヒルベルト空間、一般相対性理論、反ド・シッター時空、ラムダ項、D-brane,anti-D-brane,コンフォーマル場、
   ホログラム空間、ストリング理論、収率による商代数、ニュートンポテンシャルエネルギー、剛体力学、
   統計力学、熱力学、量子スピン、半導体、超伝導、ホイーストン・ブリッチ回路、非可換確率論、Connes理論、
   これら、演算子代数を形成している、微分・積分作用素が、ヒルベルト空間に存在している多様体の特質を
   全面に押し出して、いろいろな多様体と関数そして、群論を形作っている。


$$ \begin{vmatrix} D^m & dx \\
dx & \sigma^m \end{vmatrix} \begin{vmatrix} \cos \theta & - \sin \theta \\
\sin \theta & \cos \theta \end{vmatrix} \begin{vmatrix} x \\
y \end{vmatrix} = \begin{pmatrix} 1 & 0 \\ 
0 & - 1 \end{pmatrix}^{1 \over 2} $$
$$ \sigma^m \begin{bmatrix} \delta(x) & - 1 \\
1 & \epsilon(x) \end{bmatrix}^{1 \over 2} = \begin{pmatrix} i & 0 \\
0 & - i \end{pmatrix} $$
$$ (D^m,dx)\cdot (\cos \theta, \sin \theta) \times (dx, \partial^m)\cdot (\cos \theta, \sin \theta) $$

$$ \begin{pmatrix} 1 && 0 \\
   0 && i \end{pmatrix} \beta(x,\theta) \begin{pmatrix} x \\
   y \end{pmatrix} $$
   $$ = \begin{pmatrix} i && 0 \\
   0 && - 1 \end{pmatrix} $$
   $$ l = {\sqrt{\hbar G \over c^3}}, \sigma^m \cdot \begin{pmatrix} \delta(x) && - 1 \\
   1 && \epsilon (x) \end{pmatrix}^{1 \over 2} = \begin{pmatrix} i && 0 \\
   0 && - i \end{pmatrix} $$
   String theorem also imaginary equation of pole with the fields of manifold.
   $$ \left({\partial \over \partial \tau}f(x,y,z)\right)^{3'} = A^{\mu\nu} $$
   $$ {d \over dt}g_{ij}(t) = - 2 R_{ij} $$
   Rich equation is all of top formula in integrate theorem.
      $$ = (D^m,dx) \cdot (\cos \theta, \sin \theta) \times (dx, \partial^m)\cdot (\cos \theta, \sin \theta) $$
      Dalanverle equation equals with abel manifold.




\end{document}
